\documentclass{article}
\usepackage[utf8]{inputenc}
\usepackage[spanish,es-tabla]{babel}
\usepackage{graphicx, wrapfig}
\usepackage{anyfontsize}
\usepackage[spanish]{babel}
\usepackage{geometry}
\usepackage{fancyhdr}
\usepackage{setspace}
\usepackage{hyperref}
\usepackage{array}
\usepackage{float}
\usepackage{listings}
\usepackage{xcolor}

% fuente del título
\usepackage[T1]{fontenc}


\colorlet{punct}{red!60!black}

\colorlet{numb}{magenta!60!black}

\usepackage{xcolor}
\usepackage{vmargin}
\usepackage{pdfpages}
\usepackage{svg}
\usepackage{booktabs}
\usepackage{amsmath} % you need amsmath as the demo includes a use of \eqref
\usepackage{cite}
\usepackage[spanish]{babel}

\definecolor{azure}{rgb}{0.0, 0.5, 1.0}
\definecolor{pblue}{rgb}{0.13,0.13,1}
\definecolor{pgreen}{rgb}{0,0.5,0}
\definecolor{pred}{rgb}{0.9,0,0}
\definecolor{pgrey}{rgb}{0.46,0.45,0.48}
\definecolor{background}{HTML}{EEEEEE}
\definecolor{delim}{RGB}{20,105,176}
\hypersetup{
    colorlinks=true,
    linkcolor=azure,
    filecolor=magenta,      
    urlcolor=azure,
    citecolor=azure,
    hypertexnames=true,
}

\lstset{language=Java,
  showspaces=false,
  showtabs=false,
  breaklines=true,
  showstringspaces=false,
  breakatwhitespace=true,
  commentstyle=\color{pgreen},
  keywordstyle=\color{pblue},
  stringstyle=\color{pred},
  basicstyle=\ttfamily,
  moredelim=[il][\textcolor{pgrey}]{},
  moredelim=[is][\textcolor{pgrey}]{\%\%}{\%\%}
}
\lstdefinelanguage{json}{
    basicstyle=\normalfont\ttfamily,
    numbers=left,
    numberstyle=\scriptsize,
    stepnumber=1,
    numbersep=8pt,
    showstringspaces=false,
    breaklines=true,
    frame=lines,
    backgroundcolor=\color{background},
    literate=
     *{0}{{{\color{numb}0}}}{1}
      {1}{{{\color{numb}1}}}{1}
      {2}{{{\color{numb}2}}}{1}
      {3}{{{\color{numb}3}}}{1}
      {4}{{{\color{numb}4}}}{1}
      {5}{{{\color{numb}5}}}{1}
      {6}{{{\color{numb}6}}}{1}
      {7}{{{\color{numb}7}}}{1}
      {8}{{{\color{numb}8}}}{1}
      {9}{{{\color{numb}9}}}{1}
      {:}{{{\color{punct}{:}}}}{1}
      {,}{{{\color{punct}{,}}}}{1}
      {\{}{{{\color{delim}{\{}}}}{1}
      {\}}{{{\color{delim}{\}}}}}{1}
      {[}{{{\color{delim}{[}}}}{1}
      {]}{{{\color{delim}{]}}}}{1},
}

\setpapersize{A4}
\setmargins{3cm}        % margen izquierdo
{1.5cm}                 % margen superior
{15cm}                  % anchura del texto
{22.75cm}               % altura del texto
{1cm}                   % altura de los encabezados
{1cm}                   % espacio entre el texto y los encabezados
{0pt}                   % altura del pie de página
{1cm}                   % espacio entre el texto y el pie de página



\begin{document}

\begin{titlepage}
{\fontfamily{qag}\selectfont

\begin{center}
\vspace*{-2.5cm}

\includegraphics[width=\linewidth]{images/logoUC.pdf}

\vspace{1cm}

\textbf{\Huge\emph{Facultad \\ de\\ Ciencias}}

\vspace{2cm}

{\fontsize{20}{24}\selectfont \textbf{TÍTULO EN ESPAÑOL EN MAYÚSCULAS}}

\vspace{\baselineskip}

{\fontsize{19}{24}\selectfont (Título en inglés en minúsculas entre paréntesis)}

\vspace{1.5cm}

{\fontsize{14}{17}\selectfont Trabajo de Fin de Grado\\
para acceder al}

\vspace{0.5cm}

{\fontsize{17}{20}\selectfont \textbf{GRADO EN <RELLENAR>}}

\vfill

\begin{flushright}
\fontsize{17}{20}\selectfont
Autor: Nombre autor\\
Director: Nombre director\\
Co-director: Nombre codirector (opcional)\\
Mes - Año
\end{flushright}
\end{center}
}
\end{titlepage}


% Acknowledgements Page (Optional)                              
\newpage
\begin{center}
{\bf \Huge Agradecimientos}
\end{center}
\vspace{1cm}
\setlength{\baselineskip}{0.8cm}
\onehalfspacing
\begin{flushleft}

Me gustaría agradecer...


\end{flushleft}

\newpage

\thispagestyle{empty}

\begin{center}
{\bf \Huge Resumen}
\end{center}
\onehalfspacing



Resumen en Español



   




\begin{center}
{\bf \Large Palabras clave}
\end{center}
% Insertar pablabras clave  
Palabra clave 1, palabra clave 2, ...

\newpage

\thispagestyle{empty}

\begin{center}
{\bf \Huge Abstract}
\end{center}

English abstract ...


\begin{center}
{\bf \Large Keywords}
\end{center}
Keyword 1, Keyword 2, ...

\newpage


%\setcounter{page}{1}



\tableofcontents
\newpage
\listoffigures
\newpage
\listoftables
\newpage
\newpage

\section{Introducción}

Esta plantilla se ha desarrollado con el fin de facilitar la realización de la memoria del Trabajo Fin de Grado a los futuros egresados de la Facultad de Ciencias.

La siguiente sección muestra las funcionalidades básicas de latex que se aconsejan utilizar durante la escritura del manuscrito.

Si bien se puede modificar el estilo de esta plantilla como se desee, el objetivo de la misma es justo el contrario: que los estudiantes se centren en depurar el contenido de sus memorias, y no pierdan el tiempo con aspectos de presentación de las mismas si así lo desean.
\newpage

\section{Instrucciones}

\newpage


\newpage
\bibliographystyle{IEEEtran}
\bibliography{bibliografia.bib}


\end{document}
